\documentclass[10pt,a4paper]{amsart}
\usepackage[margin=19mm]{geometry}
\usepackage{amsmath,amssymb,mathtools}
\usepackage[scr=rsfs,cal=euler]{mathalfa}
\usepackage{xfrac}
\usepackage[unicode,hidelinks,bookmarks=false]{hyperref}
\newcommand{\ie}{i.e.\ }
\newcommand{\cf}{cf.\ }
\newcommand{\resp}{respectively\ }
\newcommand{\Subsection}{\subsection}
\providecommand{\nobreakdash}{\nobreak\hbox{-}}
\setlength{\emergencystretch}{2em}
\multlinegap0pt
\usepackage{amssymb}
\usepackage{mathtools}
\usepackage{xcolor}
\newtheorem{lemma}{Lemma}
\newcommand{\D}[1]{\mathcal{D}_{#1}}
\newcommand{\E}{\boldsymbol{E}}
\newcommand{\F}{\mathcal{F}}
\newcommand{\GEP}{\boldsymbol{E}^{\mathrm{EP}}}
\newcommand{\W}[1]{\mathcal{W}_{#1}}
\renewcommand{\a}{\boldsymbol{a}}
\newcommand{\err}[1]{\mathrm{err}(#1, S)}
\newcommand{\norm}[1]{\left\lVert#1\right\rVert}
\newcommand{\ones}[1]{\boldsymbol{1}_{#1}}
\renewcommand{\u}{\boldsymbol{u}}
\renewcommand{\v}{\boldsymbol{v}}
\title{Probabilistic Gradient Coding via Structure-Preserving Sparsification}
\author{Yuxin Jiang, Wenqin Zhang, Lele Wang}
\date{Source: arXiv:2604.10374v2 (CC BY 4.0)}
\begin{document}
\maketitle

\noindent\emph{Source and licence.} Probabilistic Gradient Coding via Structure-Preserving Sparsification. Version arXiv:2604.10374v2 (CC BY 4.0), distributed by arXiv under the Creative Commons Attribution 4.0 International licence, http://creativecommons.org/licenses/by/4.0/, as recorded in the arXiv licence metadata for this exact version and shown on the abstract page https://arxiv.org/abs/2604.10374v2. Full licence notice: \texttt{licenses/arxiv-2604.10374-CC-BY-4.0.txt}.\par\smallskip

\noindent\textit{Editorial scope.} Counted unit: the article's lemma lem:biadj-matrix (source0.tex lines 1226--1234). The article's complete original proof of it is delivered with it (source0.tex lines 1929--1960, one \texttt{proof} environment of 32 lines, which the article places away from the statement). No mathematics is written, repaired or re-derived: the statement and the proof are the article's own lines, in the article's own notation, and no retained line is substituted or re-flowed.  Retained uncounted as the source-local context the proof needs: lines 275--289.  Nothing was omitted from any retained span, so the package carries no omission record.  No retained line contains a citation, so the wrapper carries no bibliography and no bibliography entry.  No retained line contains a figure environment, an image-inclusion command or drawing code, so no image is delivered and the package declares no asset dependency.  Editorial formatting only, in addition: an AMS \texttt{amsart} preamble that restates the article's own declarations and macros used by the retained lines, namely the environments \texttt{lemma} on separate counters, together with the standard packages those lines need.  The retained environments print in this document as Lemma 1 in order of appearance, whereas the article prints the same items with the numbering of its own full document.  The complete native source of the article as distributed in the arXiv e-print for this exact version is bundled unchanged as \texttt{sources/198268/source0.tex}.
An $(N, d)$-BEG-GC corresponds to a $d$-regular bipartite expander graph $G = (\W{}\cup \D{},\mathcal{E})$, where $|\W{}| = |\D{}| = N$. Let $\E^{\mathrm{BEG}}$ be the biadjacency matrix of $G$. The normalized biadjacency matrix $\bar\E^{\mathrm{BEG}} \triangleq \tfrac{1}{d} \E^{\mathrm{BEG}}$ is the encoding matrix for the BEG-GC. Fix a non-straggler set $\F \subset [N]$ and $|\F| = N-S$. Its decoding vector takes the form
\begin{equation}
    \v_\F^{\mathrm{BEG}} = \ones{N}+\a_{\F}
    , \text{ where}\quad
    \a_\F = 
    \begin{cases}
        -1, \quad i\not \in \F\\
        \frac{S}{N-S} \quad i\in \F.
    \end{cases}
\label{eq:decoding_vec}
\end{equation} Its error is upper bounded as
\begin{equation*}
%\label{eq:expander-code-error}
    \err{\bar\E^{\mathrm{BEG}}} \leq \tfrac{1}{N}\Big(\tfrac{\sigma_2(\bar\E^{\mathrm{BEG}})}{d}\Big)^2\tfrac{NS}{N-S}.
\end{equation*}%

\begin{lemma}
\label{lem:biadj-matrix}
For the encoding matrix $\GEP_\varepsilon$ of an $(N, c, d, \varepsilon)$ expansion-preserving gradient code, we conclude that
\begin{enumerate}
\renewcommand{\labelenumi}{\roman{enumi})}
    \item $(\ones{N}, \ones{N}, d)\in \{(\u_i(\GEP_\varepsilon), \v_i(\GEP_\varepsilon), \sigma_i(\GEP_\varepsilon)\}_{i=1}^n$.
    \item $\sigma_2(\GEP_\varepsilon)< d$.
\end{enumerate}
\end{lemma}

\begin{proof}
By construction,
\[
\ones{N}^\top \a_\F
= (N-S)\left(-\tfrac{S}{N-S}\right)+S\cdot 1
= -S+S
=0.
\]
Hence $\a_\F\perp\ones{N}$. By Lemma \ref{lem:biadj-matrix}, we have $\a_\F\perp\v_1(\GEP)$. Because
$\{\v_1(\GEP),\dots,\v_N(\GEP)\}$ is an orthonormal basis, this implies
\[
\a_\F\in \mathrm{span}\{\v_2(\GEP),\dots,\v_N(\GEP)\},
\]
so there exist $\alpha_2,\dots,\alpha_N$ with
$\a_\F=\sum_{i=2}^N \alpha_i\v_i(\GEP)$.

From the definition of $\a_\F$ (see Eq.\eqref{eq:decoding_vec}), we have
\[
\norm{\a_\F}_2^2
= (N-S)\left(\tfrac{S}{N-S}\right)^{\!2}
+ S\cdot 1^2
= \tfrac{S^2}{N-S}+S
= \tfrac{NS}{N-S}.
\]

Since $\{\v_i(\GEP)\}_{i=1}^N$ is orthonormal and $\a_\F\perp\v_1(\GEP)$, we have
$\norm{\a_\F}_2^2
= \textstyle\sum_{i=2}^N \bigl|\alpha_i\bigr|^2.$
Therefore, it follows that
$\sum_{i=2}^N \alpha_i^2=\tfrac{NS}{N-S}$. Thus,
$\norm{\a_\F}_2=\sqrt{\tfrac{NS}{N-S}}$.
\end{proof}

\end{document}
